\scanpage{210}
\mspage{159}

Démonstration immédiate - \unsure{ça|Ça} devrait plutôt remonter au § 1.

\uline{Corollaire 1} Soit $Q$ un carré \uline{cartésien} (cf page 158) avec $i'$ une \uline{imm. ouverte} (donc $i$ aussi), et $q$ \inserted{$W$-}\uline{lisse} (donc $p$ aussi). On suppose que $W$ satisfait l'axiome $W(5)$ de connexité. Soit $T'=X'\setminus U'$, $T=X\setminus U$, supposons que $T'\unsure{-|\longrightarrow}T$ soit un isomorphisme. Alors $Q$ est cocartésien.

% slanted
\margin{Donc il suffit $q$ $W_0$-lisse, où $W_0=\{f\in\operatorname{Fl}\mathrm{Cat}\mid\unsure{\pi_0|\Pi_0}(f)\text{ bijectif}\}$}

Il suffit de prouver la condition d)\inserted{ de prop. 15}. Considérons le carré

\diagram{\begin{tikzcd}[ampersand replacement=\&]
\& U_{/t}\arrow[r,hook,"i_{/t}"]\arrow[d,"p_{/t}"'] \& X_{/t}\arrow[d,"q_{/t}"]\\
(\unsure{\alpha|a|*})\& U'_{/t'} \arrow[r,hook,"i'_{/t'}"] \& X'_{/t'} ,
\end{tikzcd}}

je dis qu'il est \uline{cartésien}. Cela résulte de \deleted{la} seule hypothèse que $Q$ soit cartésien, sans hyp. de lissité sur $q$ ni que $i'$ soit une imm. ouverte, et se voit p. ex. en regardant le \deleted{cube} diagramme cubique :

\diagram{\begin{tikzcd}[ampersand replacement=\&,row sep=small,column sep=small]
\& U_{/t}\arrow[rr]\arrow[dl]\arrow[dd,crossing over] \& \& X_{/t}\arrow[dl]\arrow[dd]\\
U'_{/t'}\arrow[rr,crossing over]\arrow[dd] \& \& X'_{/t'}\arrow[dd]\arrow[dl,phantom,"\deleted{\unsure{\swarrow}}" description] \&\\
\& U\arrow[dl]\arrow[rr,dash,crossing over] \&\& \unsure{X'|X}\arrow[dl]\\
U'\arrow[rr,dash] \& \& X'\arrow[ul,phantom,"\deleted{\unsure{\swarrow}}" description] \&
\end{tikzcd}}

\margin{NB ici $t$ est un élément quelc. de $\operatorname{Ob}X$, $t'$ son image dans $X'$}

Dans ce diagramme, les faces \unsure{antérieures} et \unsure{postérieures} du cube \unsure{est|sont} \unsure{cartésien|cartésiens|cartésiennes}, par définition de $U_{/t}$ et \deleted{\unsure{de}} $U'_{/t'}$
% archive sheet number: 210
